\documentclass[11pt]{article}
\usepackage[margin=2.5cm]{geometry}
\usepackage{amsmath,amssymb,amsthm}
\usepackage{mathtools}
\usepackage{hyperref}

\newtheorem{theorem}{Theorem}
\newtheorem{lemma}[theorem]{Lemma}
\newtheorem{proposition}[theorem]{Proposition}
\newtheorem{corollary}[theorem]{Corollary}
\newtheorem{definition}[theorem]{Definition}
\newtheorem{remark}[theorem]{Remark}
\newtheorem{conjecture}[theorem]{Conjecture}

\newcommand{\R}{\mathbb{R}}
\newcommand{\Q}{\mathbb{Q}}
\newcommand{\Z}{\mathbb{Z}}
\newcommand{\N}{\mathbb{N}}
\newcommand{\RTCRN}{R_{\mathrm{RTCRN}}}
\newcommand{\prodpoly}{\mathrm{prod}}
\newcommand{\degrpoly}{\mathrm{degr}}

\title{Zero-Init Non-Collapse for Bounded CRN-Shape PIVPs}
\author{Xiang Huang \and Zinan Huang}
\date{2026-04-20}

\begin{document}
\maketitle

\begin{abstract}
We prove the following structural theorem about chemical reaction network (CRN)
computable numbers. Let $y' = p(y)$, $y(0) = 0$, be a $d$-dimensional polynomial
initial value problem (PIVP) admitting a \emph{CRN-shape} decomposition
$p_i(y) = \prodpoly_i(y) - \degrpoly_i(y)\cdot y_i$ with
$\prodpoly_i, \degrpoly_i$ having non-negative rational coefficients. Suppose
the trajectory is bounded on $[0, \infty)$. Then any species $y_i$ that is
strictly positive at some finite time $t_0 \ge 0$ admits a positive constant
$c > 0$ such that $y_i(t) \ge c$ on a cofinal subset of $[0, \infty)$. In
particular, $y_i(t) \not\to 0$. Consequently, if the designated output of a
bounded zero-init CRN-shape PIVP converges to $0$, then the output trajectory
is identically $0$: the constant $0$ is not non-trivially CRN-computable from
zero init. The result has been formalized in Lean~4 / Mathlib with zero custom
axioms.
\end{abstract}

%---------------------------------------------------------------------
\section{Background}
%---------------------------------------------------------------------

\subsection{CRN computability}

A \emph{polynomial initial value problem} (PIVP) of dimension~$d$ is
\begin{equation}\label{eq:pivp}
  y'(t) = p(y(t)), \qquad y(0) = y_0 \in \Q^d,
\end{equation}
where $p\colon \R^d \to \R^d$ is a polynomial vector field with rational
coefficients. PIVPs are the phase-space model for the general-purpose analog
computer (GPAC). A real number $\alpha$ is \emph{real-time CRN-computable}
($\alpha \in \RTCRN$) if it can be approximated at exponential rate by a
designated output coordinate of a bounded PIVP implementable as a chemical
reaction network with mass-action kinetics~\cite{RTCRN1,RTCRN2,LPP}.

The CRN constraint imposes strong structural requirements. Concentrations are
non-negative, reaction rates are non-negative, and any degradation of a
species $X_i$ must be catalyzed by $X_i$ itself: you can only destroy what is
already present. Abstractly, the vector field admits a decomposition
\begin{equation}\label{eq:crn-shape}
  p_i(y) = \prodpoly_i(y) - \degrpoly_i(y) \cdot y_i,
\end{equation}
where $\prodpoly_i$ and $\degrpoly_i$ are polynomials with non-negative
rational coefficients. The factor of $y_i$ in the degradation term is the
mass-action constraint: every reaction that consumes $X_i$ requires $X_i$ on
its left-hand side.

\begin{definition}[CRN-shape decomposition]\label{def:pcd}
A \emph{\mbox{PolyCRNDecomposition}} of a PIVP~\eqref{eq:pivp} is a choice of
polynomials $\prodpoly_i, \degrpoly_i \in \Q_{\ge 0}[y_1, \ldots, y_d]$
satisfying~\eqref{eq:crn-shape} for every $i$, together with
$y_0 \in \Q_{\ge 0}^d$.
\end{definition}

\subsection{Zero initialization}

CRN species start from zero concentration. There is no laboratory operation
that \emph{preloads} an arbitrary rational value into a molecular species: one
starts by filling the flask, and then the reactions run. Zero initialization
is therefore the physically canonical assumption.

\begin{definition}[Zero init]
A PIVP~\eqref{eq:pivp} is \emph{zero-initialized} if $y_0 = \mathbf{0}$.
\end{definition}

Zero init is rigid: it rules out decay as a computational mechanism. The
simplest contrast is the linear scalar system $x'=-x$, with $x(0) = 1$:
$x(t) = e^{-t} \to 0$. In the zero-init CRN-shape setting, this decay is
impossible. The field at $x = 0$ reduces to $\prodpoly(0)$, which is the
constant coefficient of $\prodpoly$; if this is positive, $x$ grows, and if it
is zero, $x$ never moves. Either way, $x$ cannot decay to a limit of $0$
non-trivially. This paper formalizes the $d$-dimensional version of that
observation.

\subsection{Root species and dependency}

A species $X_i$ is a \emph{root} if $(\prodpoly_i)_0 > 0$, that is, the
production polynomial has a strictly positive constant term. Equivalently,
there is an $\emptyset \to X_i$ reaction that synthesizes $X_i$ from nothing.
A root species can grow from zero without external input.

Non-root species depend on upstream species. We formalize this via a
graph-theoretic \emph{root-reachability} predicate: species $i$ is
root-reachable if either (i) $i$ is a root, or (ii) $\prodpoly_i$ contains a
positive-coefficient monomial $\prod_j y_j^{\sigma_j}$ every one of whose
active feeders ($j$ with $\sigma_j > 0$) is already root-reachable. This is a
least-fixed-point definition: the root-reachable set is built up inductively
from the roots.

%---------------------------------------------------------------------
\section{Main Theorem}
%---------------------------------------------------------------------

\begin{theorem}[Zero-init non-collapse]\label{thm:main}
Let $(P, \mathrm{pcd})$ be a $d$-dimensional PIVP equipped with a CRN-shape
decomposition (Definition~\ref{def:pcd}), and assume $P$ is zero-initialized.
Let $y \colon [0, \infty) \to \R^d$ be a solution of $P$ whose trajectory is
bounded:
\[
  \exists M > 0 : \; \|y(t)\| \le M \quad \text{for all } t \ge 0.
\]
If $y_i(t_0) > 0$ for some $i \in \{1, \ldots, d\}$ and $t_0 \ge 0$, then
there exists $c > 0$ such that
\[
  \forall T \ge 0,\ \exists t \ge T : \; y_i(t) \ge c.
\]
In particular, $\liminf_{t \to \infty} y_i(t) > 0$, and no limit of
$y_i(t)$ as $t \to \infty$ (if it exists) can be $0$.
\end{theorem}

The cofinal form of the conclusion is what we actually prove; it is
equivalent to $\liminf y_i \ge c$ by standard reasoning.

\begin{corollary}[Trivial-zero theorem]\label{cor:trivial-zero}
Under the hypotheses of Theorem~\ref{thm:main}, if the designated output
coordinate $y_{i^\ast}$ satisfies $y_{i^\ast}(t) \to 0$ as $t \to \infty$,
then $y_{i^\ast}(t) = 0$ for all $t \ge 0$.
\end{corollary}

\begin{proof}
Suppose $y_{i^\ast}(t^\star) \ne 0$ for some $t^\star \ge 0$. Non-negativity
of the trajectory (Proposition~\ref{prop:nonneg} below) gives
$y_{i^\ast}(t^\star) > 0$. Theorem~\ref{thm:main} provides $c > 0$ such that
$y_{i^\ast}(t) \ge c$ cofinally, contradicting $y_{i^\ast}(t) \to 0$.
\end{proof}

\begin{remark}
Corollary~\ref{cor:trivial-zero} says the constant $0$ is not a
non-trivially CRN-computable number from zero init: if $0$ is a limit of the
output, the output is literally $0$ throughout the experiment. Computing $0$
is only possible by never doing anything.
\end{remark}

%---------------------------------------------------------------------
\section{Proof}
%---------------------------------------------------------------------

The proof is modular, in three steps.

\subsection{Step 1 --- Non-negativity}

\begin{proposition}[CRN non-negativity]\label{prop:nonneg}
Under the hypotheses of Theorem~\ref{thm:main}, $y_i(t) \ge 0$ for every
$i$ and every $t \ge 0$.
\end{proposition}

\begin{proof}[Proof sketch]
This is the standard mass-action invariance. Whenever $y_i = 0$, the field
reduces to $p_i = \prodpoly_i(y) \ge 0$ by non-negativity of the coefficients
of $\prodpoly_i$ and of the remaining coordinates $y_j \ge 0$. A
squared-negative-mass Lyapunov functional $V(t) := \sum_i (y_i(t))_-^2$
satisfies $V(0) = 0$ and $V' \le L \cdot V$ for a Lipschitz constant $L$ on
the bounded trajectory ball, so Gr\"onwall gives $V \equiv 0$. Full details
appear in the Mathlib-style proof \texttt{pivp\_solution\_nonneg} in our Lean
formalization, invoked together with a local Lipschitz estimate for the
polynomial vector field.
\end{proof}

\subsection{Step 2 --- Positive lower bound for root species}

The structural heart of the proof lives here. We show that any root species
admits a positive asymptotic lower bound via a scalar Gr\"onwall argument.

\begin{proposition}[Root species lower bound]\label{prop:step2}
Under the hypotheses of Theorem~\ref{thm:main}, let $r$ be a root species, so
$c_r := (\prodpoly_r)_0 > 0$. Then there exists $c > 0$ such that
$y_r(t) \ge c$ holds on a cofinal subset of $[0, \infty)$.
\end{proposition}

\begin{proof}
On the bounded trajectory ball $B_M := \{y \in \R^d_{\ge 0} : \|y\| \le M\}$,
Step~1 gives $y(t) \in B_M$ for all $t \ge 0$. Two algebraic facts about
polynomials with non-negative coefficients on the non-negative orthant do the
work.

\smallskip
\noindent\emph{(L) Constant-coefficient lower bound.} For any polynomial
$q \in \Q_{\ge 0}[y_1, \ldots, y_d]$ and any $y \in \R^d_{\ge 0}$,
\[
  q_0 \le q(y).
\]
This holds because $q(y) = q_0 + \sum_{\sigma \ne 0} q_\sigma \prod_j y_j^{\sigma_j}$
and every summand on the right is non-negative.

\smallskip
\noindent\emph{(U) Uniform upper bound on $B_M$.} For any polynomial
$q \in \Q_{\ge 0}[y_1, \ldots, y_d]$, the quantity
$U_q(M) := \sum_{\sigma \in \mathrm{supp}(q)} q_\sigma \, M^{|\sigma|}$
(where $|\sigma| = \sum_j \sigma_j$) is a uniform upper bound:
$q(y) \le U_q(M)$ for every $y \in B_M$.

\smallskip
Applied to $\prodpoly_r$ and $\degrpoly_r$ at $y = y(t)$, (L) and (U) give
\[
  y_r'(t) = \prodpoly_r(y(t)) - \degrpoly_r(y(t)) \cdot y_r(t)
         \ge c_r - D_r \cdot y_r(t),
\]
where $D_r := U_{\degrpoly_r}(M) \ge 0$. Notice how the CRN shape is used:
the degradation term is linear in $y_r$, which is what makes the right-hand
side a scalar affine inequality in $y_r$.

Set $K := D_r + 1$ (strictly positive even when $D_r = 0$), and define the
target level $\alpha := c_r / K > 0$. Consider the deficit
$f(t) := \alpha - y_r(t)$. From the above,
\[
  f'(t) = -y_r'(t) \le -c_r + D_r \cdot y_r(t)
        = -c_r + D_r (\alpha - f(t)) = -\frac{c_r}{K} - D_r \cdot f(t)
        = -\alpha - D_r \cdot f(t).
\]
Zero init gives $f(0) = \alpha$. Scalar Gr\"onwall with slope $-D_r$ and
drift $-\alpha$ yields, for $t \ge 0$,
\[
  f(t) \le \alpha \cdot e^{-D_r t} - \frac{\alpha}{D_r}\bigl(1 - e^{-D_r t}\bigr)
         \xrightarrow{\,t \to \infty\,} -\frac{\alpha}{D_r}\ \text{(if $D_r > 0$)},
\]
with the obvious modification $f(t) \le \alpha - \alpha t \to -\infty$ when
$D_r = 0$. In either case, for $t$ beyond an explicit threshold $t_\star$
(depending only on $c_r$, $D_r$, and $M$) we have $f(t) \le \alpha/2$, i.e.,
$y_r(t) \ge \alpha/2$. This is the required positive lower bound, holding
for all $t \ge t_\star$ (hence cofinally).
\end{proof}

Proposition~\ref{prop:step2} is actually proved in two forms in the Lean
file: the cofinal form used directly above, and the stronger
\emph{eventual} form asserting $y_r(t) \ge c$ for \emph{all} $t \ge T$ and
some $T$. The latter is what feeds into Step~3.

\subsection{Step 3 --- Propagation to root-reachable species}

Step 2 covers root species only. A non-root species $i$ that ever becomes
positive must be \emph{fed}, directly or transitively, by a root. Step 3 has
two components: a graph-theoretic statement and an analytic propagation.

\subsubsection{Root-reachability propagates lower bounds}

\begin{proposition}[Analytic propagation]\label{prop:analytic-prop}
If every root species $r$ admits an eventual positive lower bound, then so
does every root-reachable species~$i$.
\end{proposition}

The proof is an induction on the inductive definition of root-reachability.
The root case is Step~2. For the step case, if $\prodpoly_i$ contains a
positive-coefficient monomial
$m(y) = q_\sigma \prod_{j : \sigma_j > 0} y_j^{\sigma_j}$ every active feeder
of which has an eventual lower bound, then $m(y(t))$ itself has a positive
eventual lower bound (product of positive quantities is positive). Since
every other monomial of $\prodpoly_i$ is non-negative, $\prodpoly_i(y(t))$
inherits that positive lower bound, and the same scalar Gr\"onwall argument
as in Step~2 applies to $y_i$.

\subsubsection{Graph traversal: ever-positive implies root-reachable}

The non-trivial content is the converse direction of the graph: if $i$ is
ever positive on the trajectory, then $i$ is root-reachable. Phrased in the
contrapositive: if $i$ is not root-reachable, then $y_i(t) = 0$ for every
$t \ge 0$.

\begin{proposition}[Graph traversal]\label{prop:graph-trav}
Let $R \subseteq \{1, \ldots, d\}$ be the set of root-reachable species. For
every $t \ge 0$ and every $j \notin R$, $y_j(t) = 0$.
\end{proposition}

\begin{proof}
Let $N := \{1, \ldots, d\} \setminus R$ denote the \emph{dead} species. We
study the quadratic functional
\begin{equation}\label{eq:S-functional}
  S(t) := \sum_{j \in N} y_j(t)^2.
\end{equation}
Zero init gives $S(0) = 0$, and Proposition~\ref{prop:nonneg} combined with
the CRN-shape polynomial structure gives $S(t) \ge 0$ for all $t \ge 0$. The
proof reduces to showing $S' \le C \cdot S$ on $[0, \infty)$ for some
constant~$C \ge 0$ depending only on the polynomial structure and $M$;
scalar Gr\"onwall then forces $S \equiv 0$, and since $S$ is a sum of
squares, every $y_j$ with $j \in N$ is identically zero.

Differentiate~\eqref{eq:S-functional}:
\[
  S'(t) = 2 \sum_{j \in N} y_j(t) \cdot p_j(y(t))
        = 2 \sum_{j \in N} y_j(t) \cdot \bigl( \prodpoly_j(y(t)) -
            \degrpoly_j(y(t)) \cdot y_j(t) \bigr).
\]
The degradation contribution is $-2 \sum_{j \in N} \degrpoly_j(y(t))\, y_j(t)^2
\le 0$ (non-negative coefficients, non-negative values), so we may drop it:
\begin{equation}\label{eq:S-prime-bound}
  S'(t) \le 2 \sum_{j \in N} y_j(t) \cdot \prodpoly_j(y(t)).
\end{equation}
The crux is now a structural observation about the monomials of
$\prodpoly_j$ for $j \in N$.

\smallskip
\noindent\textbf{Claim.} For every $j \in N$ and every monomial $\sigma$ with
$(\prodpoly_j)_\sigma > 0$, there exists some $k = k(j, \sigma) \in N$ with
$\sigma_k > 0$ (i.e., at least one active feeder of the monomial is itself
dead).

\smallskip
\noindent\emph{Proof of claim.} If the claim failed, every feeder $k$ with
$\sigma_k > 0$ would be root-reachable. But then $j$ would be root-reachable
via the $\texttt{step}$ clause of the inductive definition (apply the step
clause to this monomial, using root-reachability of all its feeders),
contradicting $j \in N$. \hfill$\square$

\smallskip
Using the claim, bound each summand $y_j \cdot m(y)$ for a monomial
$m(y) = (\prodpoly_j)_\sigma \prod_\ell y_\ell^{\sigma_\ell}$ by the
``dead-feeder'' factor $y_k$:
\[
  y_j \cdot m(y) \le (\prodpoly_j)_\sigma \, M^{|\sigma|} \cdot y_j \cdot y_k
  \le (\prodpoly_j)_\sigma \, M^{|\sigma|} \cdot \tfrac{1}{2}(y_j^2 + y_k^2),
\]
where we used $y_\ell \le M$ on the remaining factors and AM--GM on the final
$y_j \cdot y_k$. Both $j \in N$ and $k \in N$, so both $y_j^2$ and $y_k^2$
are summands of $S$. Summing over all $(j, \sigma)$ with
$(\prodpoly_j)_\sigma > 0$, we obtain
\[
  S'(t) \le C \cdot S(t), \qquad t \ge 0,
\]
for an explicit finite constant~$C$. Combined with $S(0) = 0$ and
$S \ge 0$, scalar Gr\"onwall yields $S(t) \le 0 \cdot e^{C t} = 0$ for all
$t \ge 0$. Hence $y_j(t) = 0$ for every $j \in N$ and every $t \ge 0$.
\end{proof}

\begin{remark}
The $S$-functional approach bypasses a continuity argument based on the
first time a species becomes positive, which we initially considered. Since
zero init already guarantees $S(0) = 0$, and scalar Gr\"onwall directly
propagates that to $S \equiv 0$, no first-positive-time analysis is needed.
This simplification came from recognizing that the dead-species set is
itself forward-invariant, which is encoded cleanly in a single quadratic
functional.
\end{remark}

\subsubsection{Assembling the theorem}

By Proposition~\ref{prop:graph-trav}, the hypothesis $y_i(t_0) > 0$ implies
$i \in R$, i.e., $i$ is root-reachable. By Proposition~\ref{prop:analytic-prop}
(with the eventual form of Step~2 as base case), $y_i$ admits an eventual
positive lower bound, which implies the cofinal form of
Theorem~\ref{thm:main}. \hfill$\square$

%---------------------------------------------------------------------
\section{Lean 4 Formalization}
%---------------------------------------------------------------------

The proof has been mechanized in Lean~4 with Mathlib. The main file is
\begin{center}
\texttt{Ripple/Core/ZeroInitPositivity.lean} (1635 lines).
\end{center}

\subsection*{Headline theorems}

\begin{itemize}
  \item \texttt{zero\_init\_no\_collapse} (line 1590) --- Theorem~\ref{thm:main}.
  \item \texttt{zero\_of\_limit\_is\_trivial} (line 1610) --- Corollary~\ref{cor:trivial-zero}.
  \item \texttt{crn\_trajectory\_nonneg} (line 162) --- Proposition~\ref{prop:nonneg}.
  \item \texttt{noCollapse\_step2\_root\_liminf} (line 362) and
        \texttt{noCollapse\_step2\_root\_eventual} (line 782) --- Proposition~\ref{prop:step2}, in cofinal and eventual forms respectively.
  \item \texttt{noCollapse\_step3\_scc\_induction} (line 1552) ---
        Proposition~\ref{prop:analytic-prop}, via
        \texttt{rootReachable\_hasEventualLowerBound} (analytic propagation) and
  \item \texttt{everPositive\_rootReachable} (line 1157) ---
        Proposition~\ref{prop:graph-trav}, proved via the
        \mbox{$S$-functional} Gr\"onwall described in Section~3.3.2.
\end{itemize}

\subsection*{Axiom trace}

The Lean command
\begin{center}
  \texttt{\#print axioms zero\_init\_no\_collapse}
\end{center}
returns exactly \texttt{[propext, Classical.choice, Quot.sound]}, the three
standard Mathlib axioms. The formalization uses \textbf{zero custom axioms}:
every structural claim is discharged to a mechanical proof.

\subsection*{Milestones}

The proof was developed incrementally over several commits:

\begin{center}
\begin{tabular}{ll}
\hline
Commit & Milestone \\
\hline
\texttt{4ac4c00} & Initial formalization scaffold (three axiom placeholders) \\
\texttt{abe1527} & Step~2 discharged: root species Gr\"onwall lower bound \\
\texttt{9c6ab11} & Step~3 graph-traversal theorem (modular form) \\
\texttt{b1eebc6} & Rank-form descent for \texttt{everPositive\_rootReachable} \\
\texttt{c72484f} & Final axiom discharged: \texttt{everPositive\_hasFeedingMonomial} \\
\texttt{bd1f1e9} & \texttt{polyPIVP\_field\_locally\_lipschitz} promoted to theorem \\
\hline
\end{tabular}
\end{center}

After \texttt{c72484f} the theorem is axiom-free beyond Mathlib's three
standard axioms.

%---------------------------------------------------------------------
\section{Discussion}
%---------------------------------------------------------------------

\subsection{What the theorem rules out}

Theorem~\ref{thm:main} is a \emph{non-collapse} statement: a bounded
zero-init CRN-shape species cannot ``decay to zero'' in any limiting sense,
once it has ever become positive. Equivalently,
Corollary~\ref{cor:trivial-zero} says $0$ is computable from zero init only
trivially, by never producing any output signal.

The obvious question is whether the theorem leaves room for weaker forms of
collapse: can $y_i(t)$ approach $0$ arbitrarily closely without converging to
it? The answer is no: the cofinal conclusion gives a uniform $c > 0$ such
that $y_i(t) \ge c$ happens for arbitrarily large $t$, so
$\liminf_{t \to \infty} y_i(t) \ge c > 0$, which rules out
$\liminf y_i(t) = 0$ in particular.

\subsection{Why non-zero init is different}

Without zero init, the conclusion fails spectacularly. The one-dimensional
system $x' = -x$, $x(0) = 1$ is CRN-shape ($\prodpoly = 0$, $\degrpoly = 1$),
bounded, and satisfies $x(t) = e^{-t} \to 0$. The non-negative-init condition
alone does \emph{not} force the conclusion --- it is the zero-init condition
that forces production to dominate at every zero state.

\subsection{Connections and forward directions}

\paragraph{Positivity certificates.} The monomial-cover argument in the
claim within Proposition~\ref{prop:graph-trav} is effectively a Positivstellensatz
witness for the forward-invariance of the dead-species orthant. An interesting
open question is whether the inductive root-reachability predicate admits a
polynomial-time syntactic check (it does: it is a reachability problem on
the monomial support hypergraph) and whether more refined reachability
concepts (e.g., reachability with quantitative rate bounds) can be formalized
similarly.

\paragraph{Signomial programming.} The scalar Gr\"onwall argument in Step~2
gives an explicit bound $\alpha := c_r / (D_r + 1)$, and this is provably
tight up to the $+1$ regularizer: the asymptotic $\liminf$ of $y_r$ is
$c_r / D_r$ in the case $D_r > 0$. A tighter quantitative theory, with
$\alpha$ as a function of the polynomial data and the box size $M$, would
connect directly to signomial programming and to the bounded analog
complexity hierarchy of~\cite{BAC}.

\paragraph{Characterizing LPP-computability.} Theorem~\ref{thm:main} is a
\emph{necessary} condition for any species in a zero-init CRN to play the
role of a non-trivial output. Combined with the constructive CRN
realizations from~\cite{RTCRN1,RTCRN2,LPP}, it sharpens the characterization
of $\RTCRN$ from zero-init CRNs: the only limit points of output trajectories
that are genuinely computed (as opposed to being identically zero) are
strictly positive.

%---------------------------------------------------------------------
\begin{thebibliography}{9}

\bibitem{RTCRN1}
X.~Huang, T.H.~Klinge, J.I.~Lathrop, X.~Li, J.H.~Lutz.
\newblock Real-time computability of real numbers by chemical reaction networks.
\newblock \emph{Natural Computing}, 2018.

\bibitem{RTCRN2}
X.~Huang, T.H.~Klinge, and J.I.~Lathrop.
\newblock Real-time equivalence of chemical reaction networks and analytic
  signal transducers.
\newblock In \emph{Proc.\ DNA~25}, LNCS 11648, pp.~31--50, 2019.

\bibitem{LPP}
X.~Huang and R.~Huls.
\newblock Computing real numbers with large-population protocols.
\newblock In \emph{Proc.\ DNA~28}, 2022.

\bibitem{BAC}
X.~Huang.
\newblock Bounded analog complexity.
\newblock Submitted, 2026.

\end{thebibliography}

\end{document}
